% Source: https://artofproblemsolving.com/community/c3h1737416_geometry_problem
% Stable identity: aops-topic-1737416-q01
By the law of cosines,
\[
\cos A=\dfrac{AB^2+AC^2-BC^2}{2\cdot AB\cdot AC}=\dfrac{196+169-225}{364}=\dfrac{5}{13},\qquad \sin A=\dfrac{12}{13}.
\]
Angle $A$ is acute, so the circumcentre $X$ lies on the same side of $BC$ as $A$, and the central angle is $\angle BXC=2\angle A$:

\begin{asy}
usepackage("amsmath");
size(6cm);
defaultpen(fontsize(10pt));
pair B=(0,0),C=(15,0),A=(8.4,11.2),X=(7.5,3.125);
draw(circle(X,8.125),gray(0.55));
draw(A--B--C--cycle,linewidth(0.9));
draw(B--X--C,blue+linewidth(0.8));
dot(X);
draw(arc(X,1.0,degrees(B-X),degrees(C-X)),blue);
label("$2\angle A$",X+2.15*dir(270),blue);
label("$A$",A,N);label("$B$",B,W);label("$C$",C,E);label("$X$",X,N);
label("$14$",(A+B)/2,NW);label("$13$",(A+C)/2,NE);label("$15$",(B+C)/2,S);
label("\mbox{}",(min(currentpicture,user=true)+max(currentpicture,user=true))/2,black); // keeps the text after the figure black when LaTeX typesets the labels inline
\end{asy}

Let $\theta=180^\circ-2A$ be its supplement. Then
\[
\cos\theta=-\cos2A=1-2\cos^2A=\dfrac{119}{169},\qquad \sin\theta=\sin2A=2\cdot\dfrac{12}{13}\cdot\dfrac{5}{13}=\dfrac{120}{169}.
\]
In a right triangle with an acute angle $\theta$, the adjacent leg, opposite leg and hypotenuse are in the ratio $\cos\theta:\sin\theta:1=119:120:169$. Since $\gcd(119,120)=1$, the integer triangles with this angle are the multiples $119k,120k,169k$, and the smallest perimeter is
\[
119+120+169=\boxed{408}.
\]
